\section{Confrontation: classical cartoon, quantum framework, phenomenology}
\label{sec:confrontation}

Three independent handles now aim at the same ratios:

\begin{enumerate}
\item the classical simulation of Secs.~\ref{sec:model}--%
\ref{sec:classical-results} (cheap, dynamical, schematic color);
\item the Tier-1--4 lattice observables of Sec.~\ref{sec:hierarchy}
(rigorous, awaiting hardware at scale);
\item phenomenological baselines: statistical hadronization at
$T_{\rm ch}\approx156\MeV$ \cite{Andronic:2017pug}, coalescence
\cite{Fries:2008hs}, and the ExHIC exotic-yield systematics
\cite{ExHIC:2010gcb,ExHIC:2017smd,Zhang:2020dwn}, against data anchors such as
$\chic$ production in pp and Pb--Pb \cite{CMS:2021znk,LHCb:2021auc} and
the $\Tcc$ discovery \cite{LHCb:2021vvq}.
\end{enumerate}

The crosswalks are explicit by construction.  The classical
manifest-exotic tally ($cc\bar u\bar d$ versus $c\bar c q\bar q$) is the
same partition Tier~1 draws with conserved charges.  The classical binding
criterion maps onto Tier-4 flowed projectors, with
$r_0 \leftrightarrow \sqrt{8t_{\rm fl}}$ as the scheme dial; comparing the
$\tau_{\rm bound}$-scanned classical ratios with $t_f$-scanned Tier-1--3
counts separates formation from survival in both frameworks.
\todo{when production scans exist: table of M:B:T:P from all three
handles, normalized to the meson yield.}
